\section{Efficient Flagship, Vendor-Framed}
\sectionkicker{EFFICIENT FLAGSHIP}
\begin{wideflow}
{\Large\bfseries 効率の旗艦を、発表の範囲で読む}\par
\smallskip
{\color{SurveyMuted}9月10日のDeepSeek V4.1-Flashを、他社超えの言い切りではなく提供側の報告として受け止め、9月14日の切り替えは先の話として分ける。}\par
\end{wideflow}

9月10日、DeepSeekはV4.1-Flashを示した\cite{deepseekflash}。5520億規模の混合専門家の構えで、入力で80億、出力で160億が働き、画像の理解を素で備え、100万の文脈を持つとされる。学び直しと強化学習の積み増しが添えられる。重みの置き場はMITの許しで開かれ、APIではdeepseek-flashとして生きる。旧いV4-FlashやVision-Expは引退し、互換の経路が用意されるとされる。省く側の記憶域はHBMで4分の1、SSDで8分の1と説明され、代理人の費用に効くとされる。表や図の細目は取得の打ち切りでたどり切れておらず、採点の取り方は提供側のまとめとして扱う。給仕の対応は約束の段階であり、引擎側の時点の状態と切り分ける。

比較の言い方は慎重に扱う。提供側は複数の外部の手による試しを引きつつ、V4-Proを含む旗艦群を上回ると述べるが、試した相手の名は明かされず、独立の再現もない\cite{deepseekflash}。本号では「V4.1-FlashがV4-Proに勝つ」と言い切らず、提供側の報告として書き分ける。9月14日4時（世界時）にV4-Proへの求めをV4.1-Flashへ振り替える案内は、範囲の中の文書で予告された範囲の後の運用であり\cite{deepseekapi}、今週起きた出来事として書かない。範囲の中の料金改定（9月10日4時から有効、繁閑で半額の時間帯あり）は別に起きた事実として扱い、表の1符号あたりの値は図に埋もれて取り切れていないことを添える。V4.1-Proの出日は資料に見当たらない。

\begin{claimboundary}[比較と日付の境界]
比べた相手の名は明かされていない。9月14日の振り替えは先の運用であり、今週の出来事ではない。表の細値は取り切れていない。
\end{claimboundary}

\begin{communitynote}[X / Community observation --- context only]
9月10日6時10分（世界時）の公式発信に、独立の試しや計測の報告が重なる様子が観察される\cite{w37x-deepseek-official,w37x-deepseek-ind}。うかがえるのは導入や計測への関心の向きであり、性能や費用の根拠にはならない。
\end{communitynote}
